\cotefile{134-2}
\scanpage{184}\mspage{39}
on arrive à réconstituer le type d'homotopie via son $\infty$-groupoïde fondamental, comme expliqué dans ma lettre précédente, \deleted{\unsure{e|a}}et rappelé dans celle-ci. (Cela signifie en particulier que toutes les opérations cohomologiques possibles et imaginables sont déjà incluses dans la donnée d'un tel système de $n$-catégories .\,.\,.\,.) \deleted{Mais}

De même, les types d'homotopie plus \unsure{cannelés|canelés}, liés à certains diagrammes finis, qu'on peut associer à \deleted{\unsure{des}}certains types de stratifications \inserted{$\mathcal{P}$} d'espaces topologiques ``modérés'' $X$ disons, devraient correspondre de façon aussi parfaite aux \deleted{\unsure{$n$-}}\inserted{$(n+1)$-}catégories de $n$-champs sur $X$ qui sont loc constant sur chacune des strates de $\mathcal{P}$ (disons: qui sont subordonnés à $\mathcal{P}$). Si plus haut le \unsure{topos|topis} de la cat. dérivée loc. constante marchait bel et bien, on s'attendrait ici à récupérer le type d'homotopie mixte à partir de la \inserted{ss-}catégorie de la catégori dérivée des faisceaux abéliens sur $X$, formée des complexes qui sont à \unsure{coh.|ch.} loc. constant sur chacune des strates -- avec en plus les opérations $\overset{\mathbb{L}}{\otimes}$, $\mathbb{R}\underline{\mathrm{Hom}}$ plus, les
% archive sheet number: 184

\scanpage{185}\mspage{40}
\unsure{Cas|cas} échéant, les 4 opérations $\mathbb{R}g_!$ $\mathbb{R}g_*$, $\mathrm{L}g^*$, $g^!$ pour des inclusions $g:Z'-Z''$ de réunions localement fermées de strates.\,.\,.\,.\,.\,. (L'ennui ici est qu'on ne sait même pas encore ce que c'est qu'une catégorie triangulée, pas plus qu'on ne sait ce qu'est la version non commutative, sans doute plus simple et plus fondamentale: \uline{une} ``catégorie homotopique'', avec opérations ``fibre'' et ``cofibre'' .\,.\,.\,.)

\begin{center}\rule{0.41\linewidth}{0.4pt}\end{center}

Il est vraiment temps que je termine cette lettre-fleuve, qui devient de plus en plus vaseuse. Juste une question: Quelle est cette mirifique formule de \unsure{Bloch|Blod}-Quillen à laquelle tu fais allusion, dont je n'ai jamais entendu parler, et qui fait venir l'eau à la bouche?

Bien cordialement à toi\\
\hspace*{0.42\linewidth}Alexandre\\
\hspace*{0.52\linewidth}Grothendieck\\
% archive sheet number: 185
